\documentclass[11pt,a4paper]{article}
\usepackage{fontspec}
\setmonofont{DejaVu Sans Mono}[Scale=0.82]
\usepackage[margin=2.4cm]{geometry}
\usepackage{amsmath,amssymb}
\usepackage{booktabs}
\usepackage{xcolor}
\usepackage[colorlinks=true,linkcolor=blue,citecolor=blue,urlcolor=blue]{hyperref}

\newcommand{\tGL}{t_{\mathrm{GL}}}
\newcommand{\code}[1]{\texttt{#1}}
\newcommand{\gt}{\tilde\gamma}

\title{juliaTDGL 2: a guide to the N-band generalised\\
time-dependent Ginzburg--Landau solver}
\author{J.~A.~Robles Calder\'on\\
\small Departamento de F\'isica, Universidad Nacional de Colombia, Bogot\'a}
\date{Release 2.0 (21 October 2026)}

\begin{document}
\sloppy
\maketitle

\begin{abstract}
\noindent
Release 2.0 (\code{tdgl2/}, package \code{TDGL} 2.0.0) generalises the two-band solver of release 1.2 to any number of
bands, each with its own Kramer--Watts-Tobin (KWT) relaxation, and adds an open boundary for the self-field, an
insulating buffer, a material per region, Runge--Kutta--Chebyshev (or Legendre) stages and a rotating frame for the bands
that KWT makes stiff, and a
measuring protocol that brackets an onset by the delay of the first phase slip. With the new terms off it reproduces
release 1.2. This guide gives the model, its discretisation, the time step, the boundaries, the protocol, the optional
terms, how to run the code and what each check establishes. What is inherited unchanged from 1.2 (the lattice, the
link variables, the geometry and its editor) is described in the guide of release 1.2 (\code{docs/solver\_guide.pdf})
and only summarised here.
\end{abstract}

\tableofcontents

% ======================================================================
\section{What the solver computes}

The solver integrates the generalised time-dependent Ginzburg--Landau (gTDGL) equations of a superconductor with $N$
order parameters $\Psi_\mu$, coupled by a Josephson matrix, together with the vector potential $\mathbf A$, on a
three-dimensional box with a buffer of a few cells, under an applied field and a transport current. It was written
for the geometric superconducting diode in two-gap materials, where a sample whose shape breaks the mirror symmetry
along the direction of vortex motion carries different critical currents $I_c^\pm$ in the two directions, and the
efficiency is
\begin{equation}
  \eta = \frac{I_c^{+}-|I_c^{-}|}{I_c^{+}+|I_c^{-}|}.
\end{equation}

A run is a current ladder for each polarity. Its output is, at every current, a verdict (silent, dissipating, normal)
with the voltage, the number of vortex crossings, the band-resolved currents and the order-parameter amplitudes;
the brackets of the onsets and of $\eta$; and a machine-readable summary with the units, the code version and the
numbers that place the model in its regime (section~\ref{sec:regime}).

What release 2.0 adds over 1.2, in one line each:
\begin{itemize}
\item $N$ bands and a symmetric Josephson matrix; passive bands (condensate induced by the coupling).
\item KWT relaxation per band: the charge imbalance of Gurevich and Vinokur at the contacts, in the gauge $\Phi = 0$.
\item First-order Runge--Kutta--Chebyshev stages (by default) or Runge--Kutta--Legendre (RKL) stages for the stiff bands, and a
  rotating frame that makes the stationary
  states with charge imbalance fixed points of the step.
\item Self-field boundaries \code{:bar} (the default) and \code{:open} (free-space field by a lattice Green function).
\item An insulating buffer outside the bar of the sample and its contacts, and a sculpted sample (vacuum inside the box: a cable, a hole), with leads of its own section.
\item The onset by the delay of the first phase slip; a verdict for a sample gone normal.
\item Runs from a TOML file, materials as files, a JSON summary.
\end{itemize}

What it does not have: the scalar potential of 1.2 (its Coulomb projection, Thomas--Fermi screening and diffused
residual), the \code{Strict} numerical mode and the explicit default of 1.2. For those, and to reproduce the published
results of release 1.x, release 1.2 stays in the repository (\code{julia/}).

% ======================================================================
\section{The model}

\subsection{Free energy and units}

Lengths are in $\xi_1$, the coherence length of band 1, time in $\tGL = \xi_1^2/D_1$, fields in $H_{c2}$ of band 1 and
currents in $J_0$. Band 1 has unit stiffness and relaxation, $g_1 = \Gamma_1 = 1$. The functional is
\begin{equation}
  F = \sum_\mu \Big[ g_\mu |\mathbf D\Psi_\mu|^2 - f\,\varepsilon_\mu |\Psi_\mu|^2 + \frac{\beta_\mu}{2}|\Psi_\mu|^4 \Big]
      - \sum_{\mu<\nu} \gamma_{\mu\nu} (\Psi_\mu^*\Psi_\nu + \mathrm{c.c.}) + \kappa^2 |\nabla\times\mathbf A|^2 ,
\end{equation}
with $\mathbf D = \nabla - i\mathbf A$, $f(\mathbf r)$ the condensation mask of the geometry (between 0 and 1, or above
1 where the condensate is locally stronger), and $\varepsilon_\mu = a_\mu - T$ ($a_\mu$ the value at $T = 0$; with the
heat model $T \to \theta(\mathbf r, t)$). The sign of $\gamma_{\mu\nu}$ chooses the interband phase ($s_{++}$ for
$\gamma > 0$, $s_\pm$ for $\gamma < 0$); with three or more bands and repulsive couplings the phases can be frustrated.

A band is a \code{Band(m, \(\Gamma\), a, \(\beta\), \(\gt\))} with $g = 1/m$; a model is a \code{Model} of bands, the
matrix $\gamma$, $\kappa$, the normal conductivity $\sigma_n$ and the bath temperature $T$. \code{two\_band(...)} builds
the two-band model of release 1.x: $g = (1, 1/m_r)$, $\Gamma = (1, \Gamma_r)$, $a = (1, a_{2n})$, $\beta = 1$,
$\sigma_n = 1$, $\gt = 0$.

\subsection{Equations of motion}

For each band, in the gauge $\Phi = 0$,
\begin{equation}
  \frac{\Gamma_\mu}{\sqrt{1+\gt_\mu^2|\Psi_\mu|^2}}
  \Big[ \partial_t\Psi_\mu + \gt_\mu^2\,\Psi_\mu\,\mathrm{Re}(\Psi_\mu^*\partial_t\Psi_\mu) \Big]
  = g_\mu \mathbf D^2 \Psi_\mu + (f\varepsilon_\mu - \beta_\mu|\Psi_\mu|^2)\Psi_\mu + \sum_{\nu\neq\mu}\gamma_{\mu\nu}\Psi_\nu ,
  \label{eq:band}
\end{equation}
and for the field
\begin{equation}
  \sigma_n(\mathbf r)\,\partial_t\mathbf A = \sum_\mu \mathbf J_\mu - \kappa^2 \nabla\times\nabla\times\mathbf A ,
  \qquad \mathbf J_\mu = g_\mu\,\mathrm{Im}(\Psi_\mu^*\mathbf D\Psi_\mu), \qquad \mathbf E = -\partial_t\mathbf A .
\end{equation}
With $\gt_\mu = 0$ equation~\eqref{eq:band} is the TDGL of release 1.2. The KWT factor slows the amplitude of a band
by $\sqrt{1+\gt^2|\Psi|^2}$ relative to its phase; it is the form of Kramer and Watts-Tobin for one band, applied band
by band (the extension to several bands is a model choice). $\gt_\mu$ multiplies $|\Psi_\mu|$ in the units of the code:
a literature value $\gamma_\mu = 2\tau_{E,\mu}\Delta_{\mu,0}/\hbar$ referred to the equilibrium gap becomes
$\gt_\mu = \gamma_\mu/|\Psi_{\mu,\mathrm{eq}}|$. The material files say so with every number.

\textbf{Charge imbalance.} In the gauge $\Phi = 0$ the electrochemical potential of the condensate is the rate at
which its phase turns. Where the injected normal current converts into supercurrent (the conversion layers next to the
contacts) the phase of a stationary state turns at that rate, and KWT lengthens the layers. This is the charge
imbalance of Gurevich and Vinokur. A band-resolved branch imbalance (Schmid--Sch\"on) is not in the model.

\textbf{Conservation.} The divergence of the field equation gives $\nabla\cdot(\mathbf J_s + \sigma_n\mathbf E) = 0$,
in the continuum and on the lattice (the discrete divergence of a discrete curl vanishes), so no Poisson problem is
solved.

\subsection{Materials and the normal conductivity}

$\sigma_n$ is given, or that of a dirty BCS material with these bands (\code{\(\sigma\)n = :dirty}, \code{sigma\_dirty}):
with the coefficients of Gurevich's two-band theory, $g_\mu\Gamma_\mu/\beta_\mu \propto N_\mu D_\mu$ does not depend on the
normalisation of $\Psi_\mu$, and $\sigma_n = \sum_\mu g_\mu\Gamma_\mu/\beta_\mu / 5.79$ (one band: $u = \Gamma/\sigma_n =
5.79$, the value of Kramer and Watts-Tobin). The material files in \code{materials/} carry each parameter with its source
and say which numbers are model choices:
\code{mgb2\_dirty.toml} and \code{mgb2\_clean.toml} (MgB$_2$ estimates from the literature), and \code{mgb2\_model\_I.toml},
\code{mgb2\_model\_II.toml} (the two-band models of the earlier studies, MgB$_2$-like labels, not quantitative MgB$_2$).

\textbf{A material per region.} A shape region (or a palette entry) with \code{material = "name"} takes the bands of a
\code{[materials.name]} table of the run file; the rest of the sample is \code{[material]}. Per node the band update reads its
material's $dt/\Gamma$, $\varepsilon = a - T$, $\beta$, $\gt^2$ and Josephson row; per link the weight is
\begin{equation}
  \rho_{ab} = t_{ab}\,\frac{m_{\mathrm{ref}}}{(m_a + m_b)/2}\,\tau ,
\end{equation}
the harmonic mean of the two stiffnesses (the current crosses an interface as a series resistance) times the transparency
$t_{ab} \in (0, 1]$ of the barrier between the two materials (\code{transparency = \{ other = t \}}, \code{"material"} naming the
first), and the supercurrent, the free energy and the dissipation take the same weights, so the band equations and Amp\`ere come
from one functional. $\kappa$, $\sigma_n$ and $T$ stay the sample's, and band $\mu$ of one material meets band $\mu$ of the other
(no interband tunnelling at an interface: a one-band metal next to a two-band one is given two bands). Any number of materials:
two or three regions of the same material are the uniform run to the bit; the discrete continuity holds to $10^{-15}$ across an
interface; the proximity length into a normal or weaker region matches the lattice solution $\cosh(dx/\xi_d) = 1 + \lambda
dx^2/2g$ to $4\times10^{-5}$. A table must give every key that \code{[material]} gives ($\Gamma$, $\beta$, $\gt$, the coupling):
one left out is an error, not the package default.
Example: \code{examples/bicapa\_tunel.toml}, a tunnel-coupled bilayer. A table may also give its own \code{sigma\_n} (the
conductivity of a link is then the harmonic mean of its two nodes, a series resistance, and the field solve takes the conjugate
gradients of the insulating buffer), the rate of \code{qp\_exchange} (with the same split of $\sigma_n$ between the bands in every
material: a split that changes across an interface needs the scalar potential, which the gauge $\Phi = 0$ removes), \code{absent =
[2]} (a band absent from the material: a one-band metal next to a two-band superconductor, \code{examples/metal\_una\_banda.toml}),
and \code{heat\_D\_T}; noise and anisotropic bands work per material too. Drag with materials is refused. A table may take a
material file instead of its bands (\code{file = "mgb2\_dirty.toml"}): its bands and coupling, with the sample's $\kappa$ and $T$,
and the sample's $\sigma_n$ unless the table gives \code{sigma\_n}.

\subsection{The regime numbers}
\label{sec:regime}

Every result is reported with the numbers that place the model (\code{regime(md; L)}, also in the summary of a run):
the ratio of diffusivities $D_\mu = g_\mu/\Gamma_\mu$ of the last band to the first; the interband Josephson length over
the sample length, $\ell_J/L$; the invariant $\gamma^2/(\alpha_1\alpha_2)$; $\sigma_n/\sum_\mu g_\mu$; $u_\mu =
\Gamma_\mu/\sigma_n$; $\sigma_n$ over its dirty value; and the conversion length of the normal current at a contact,
\begin{equation}
  \lambda_c^2 = \frac{\sigma_n}{\sum_\mu \tilde\Gamma_\mu |\Psi_\mu|^2}, \qquad
  \tilde\Gamma_\mu = \frac{\Gamma_\mu}{\sqrt{1+\gt_\mu^2|\Psi_\mu|^2}} ,
\end{equation}
with $L/\lambda_c$. KWT lengthens $\lambda_c$: in the two-band notched-box model at $\kappa = 2$, from $0.80\,\xi_1$ to
$2.39\,\xi_1$ with $\gt = 10$. When $L$ is not many $\lambda_c$ the middle of the sample still carries unconverted normal
current, which moves an onset by itself; keep $L/\lambda_c \gtrsim 10$. The pair numbers exist for two bands; for any
$N$ the summary has $\lambda_c$, $u_\mu$ and the two conductivity ratios.

% ======================================================================
\section{Discretisation}

\subsection{The lattice}

As in release 1.2: a staggered Cartesian mesh with $\Psi_\mu$ on the nodes, $\mathbf A$ on the links as link variables
$U = e^{-i\mathbf A\cdot d\mathbf l}$ (so lattice gauge invariance is exact) and the field on the plaquettes; the
condensation mask $f$ on the nodes; covariant Neumann conditions on the six faces of the sample, written as ghost nodes
$\Psi_{\mathrm{face}} = U\,\Psi_{\mathrm{interior}}$. Each band is a contiguous array, and every loop over nodes or links
is free of reductions, so a run gives the same bits with any number of threads.

\subsection{The band update}

Without KWT a band is advanced by an explicit Euler step of the right-hand side $R_\mu$ of~\eqref{eq:band}, a Jacobi
sweep over the interior nodes:
\begin{equation}
  \Psi_\mu^{n+1} = \Psi_\mu^n + \frac{dt}{\Gamma_\mu} R_\mu .
\end{equation}
With KWT the left-hand side is a symmetric positive operator on $\Psi$ seen as a vector of $\mathbb R^2$, inverted per
node without dividing by $\Psi$:
\begin{equation}
  \partial_t\Psi = R' - \frac{\gt^2\,\Psi\,\mathrm{Re}(\Psi^* R')}{1+\gt^2|\Psi|^2}, \qquad
  R' = \frac{\sqrt{1+\gt^2|\Psi|^2}}{\Gamma} R .
\end{equation}
The other bands enter the Josephson term at the start of the step.

\subsection{The field update}

The curl-curl term is implicit (one solve per step), the supercurrent explicit:
\begin{equation}
  \Big(I + \frac{dt}{\sigma_n}\kappa^2 K\Big)\mathbf A^{n+1} = \mathbf A^n + \frac{dt}{\sigma_n}\mathbf J_s^n
  - \frac{dt}{\sigma_n}\kappa^2 \mathbf c ,
\end{equation}
with $K$ the Yee curl-curl and $\mathbf c$ the boundary fields. On tensor-product link ranges with $dx = dy = dz$, $K$
is diagonalised by per-axis bases (the singular vectors of the one-dimensional difference), and per wavevector
$(I + aK)^{-1} = P_L + P_T/(1 + a|k|^2)$: a direct, exact solve by dense products along each axis. Only the transverse
part of $\mathbf A$ is implicit; the longitudinal part (the conversion layers) moves with $\mathbf J_s$ alone.

\textbf{Insulating buffer.} With \code{\(\sigma\)vac} the buffer outside the bar of the sample and its contacts has
conductivity $\sigma_{\mathrm{vac}} < \sigma_n$, so all of the current crosses the sample, as in a film in vacuum with
leads at its ends. The conductivity per link makes the system non-separable; it is solved by preconditioned conjugate
gradients with the separable solve as preconditioner (tolerance \code{pcg\_tol} $= 10^{-8}$ relative to the step's right-hand
side, $10^{-11}$ before 26 October 2026; its error is then three orders below the first-order error of the step in time), started
from the previous step's solution (since 26 October 2026 by default from the linear extrapolation of the last two, with the
separable preconditioner in single precision inside a flexible conjugate gradient, \code{pcg\_fast} and \code{pcg\_f32}: 14 to 6.5
iterations and $-55\,\%$ per step on the nanowire; the state within $2\times10^{-9}$, and a ladder resumes to the bit):
consecutive steps differ by little, so it takes 12 iterations where the separable guess takes 27. A ladder starts cold,
and what continues a copy of the state (the delay confirmation) gives the solver back, so that a ladder and its
resumption keep their bits. The contrast still costs: about 8 times the cost of a step with a metal buffer.

\textbf{A sculpted sample.} A shape element with \code{vacuum = true} removes the material: no superconductor and no
normal metal ($f = 0$ is a normal metal). Its nodes leave the block: $\Psi_\mu = 0$ there, the links that touch one carry
no supercurrent ($\tau = 0$), and the band update of a node next to the vacuum reads no neighbour across that link,
\begin{equation}
  \Psi_i \leftarrow \Psi_i + \sum_{\pm,\,a} k_a\,\tau_{i,\pm a}\bigl(U_{i,\pm a}\Psi_{i\pm a} - \Psi_i\bigr) + \dots,
\end{equation}
which is the zero-flux condition of a free surface, the same one the covariant Neumann ghost gives at a face of the box.
The insulating buffer follows the sample: a link conducts ($\sigma_n$) when both its nodes are in the sample or in a lead,
and the leads take the section of the face they touch, so the current enters and leaves through the section of the sample
(a cable with leads of its own section). The self-field on the outer faces is the Biot--Savart field of that section. $J$
stays the current density: the sample carries $J$ times its meshed section. A uniform state stays uniform up to the surface
to $10^{-13}$, and the static current through a cable is $J$ times its section to $2\times10^{-3}$, none of it through the
vacuum. It needs \code{sigma\_vac}; it is not with heat or drag (their stencils do not see the surface); with
\code{qp\_exchange} the links to the vacuum are cut (below).
Shape files draw figures on any face (\code{axis = "x"} or \code{"y"}) and keep the figure with \code{outside = true}; the
shape editor draws on its three views. Example: \code{examples/cable\_k2.toml}.

% ======================================================================
\section{The time step}

\subsection{Bounds}

The step is the smallest of the stability bounds, times 0.99 and \code{dt\_scale} (0.5 by default: no mode of the
stiffest band changes sign from one step to the next). For band $\mu$,
\begin{equation}
  dt_\mu = \frac{2\Gamma_\mu}{4 g_\mu \sum_i dx_i^{-2} + h_\mu}\cdot\frac{1}{\sqrt{1+\gt_\mu^2P_\mu^2}}, \qquad
  h_\mu = 3\beta_\mu P_\mu^2 - \varepsilon_\mu ,
\end{equation}
where $P_\mu$ is the amplitude of the coupled homogeneous equilibrium (interband coupling raises it above
$\sqrt{a_\mu/\beta_\mu}$, and a passive band, $a_\mu \le 0$, has one only through the coupling). The KWT factor is the
faster phase mode. The explicit supercurrent in Amp\`ere's law gives one more bound, $2\sigma_n/\sum_\mu g_\mu a_\mu/\beta_\mu$,
which binds only when $u = \Gamma/\sigma_n$ is large ($u \gtrsim 15$); the explicit curl-curl bound
$\sigma_n\,dl/(4\kappa^2)$ disappears with the implicit solve. With the heat model, its explicit diffusion adds
$0.15\,dx^2/D_T$. The summary of a run names the bound that set the step.

\subsection{Runge--Kutta--Legendre stages}

With KWT the band bound falls as $1/\sqrt{1+\gt^2P^2}$ (by 10 with $\gt = 10$). Instead of shortening the step, the
bands are advanced by first-order RKL stages (Meyer, Balsara and Aslam, J.\ Comput.\ Phys.\ 257, 594 (2014)): the step
is that of the same model without KWT, and each band takes $s$ stages,
\begin{equation}
  Y_0 = \Psi^n,\quad Y_1 = E(Y_0),\quad Y_j = \mu_j E(Y_{j-1}) + \nu_j Y_{j-2},\quad
  \mu_j = \frac{2j-1}{j},\ \nu_j = \frac{1-j}{j},\quad \Psi^{n+1} = Y_s ,
\end{equation}
with $E$ the Euler step of $dt/q$, $q = s(s+1)/2$, so $s$ stages bridge a ratio of bounds up to $q$. The linearised band
operator is (the KWT factor, symmetric positive per node) times (the Hessian of $F$): its spectrum is real, which is what
the method needs. Stage $j$ sits at the fraction $c_j = j(j+1)/(s(s+1))$ of the step, and its links and ghosts are taken
there. \code{rkl\_margin} makes the stages cover the bound times a factor (1.1 by default with KWT): the stages are
counted at the equilibrium amplitude $P$, and an amplitude above it (in the conversion layers) can push the fastest
phase mode over the edge. \code{rkl = false} gives the Euler substeps of release 1.2.

\textbf{Chebyshev stages} (\code{rkl = "rkc"}). Damped first-order Runge--Kutta--Chebyshev (Verwer, Hundsdorfer and Sommeijer,
Numer.\ Math.\ 57, 157 (1990)): $\omega_0 = 1 + \epsilon/s^2$, $\omega_1 = T_s(\omega_0)/T_s'(\omega_0)$, and the same recursion with
$\mu_j = 2\omega_0T_{j-1}/T_j$, $\nu_j = -T_{j-2}/T_j$ (all at $\omega_0$) and $E$ the Euler step of $dt\,\omega_1/\omega_0$. Its
stability polynomial is $T_s(\omega_0 + \omega_1 z)/T_s(\omega_0)$, so $s$ stages bridge $0.762\,s^2$ explicit steps with
$\epsilon = 0.5$ against $s(s+1)/2$. The damping is \code{rkc\_damping} $= 0.5$, not the textbook 0.05: with 0.05 the weakly damped
extrema of the polynomial let the splitting with the field and the other band grow a mode, and the notched KWT box diverged
within 100 steps. A band takes RKC only where it saves a stage, so a run without KWT keeps the bits of RKL. On the notched box
of \#4 with KWT it is $-5\,\%$ per step at $\gt = 10$ (stages $(3,9)\to(3,8)$) and $-10\,\%$ at $\gt = 30$ ($(6,15)\to(5,13)$),
with the same brackets, the same delay of the first slip (0.05\,\%) and the same energy balance; \code{summary.json} writes the
integrator of each band. It is the default since 26 October 2026 (\code{rkl = true} gives the RKL stages of the runs made before).
\textbf{The defaults a run folder keeps.} A run writes the scheme defaults it started with into its folder
(\code{scheme\_defaults.toml}); a folder that holds a ladder without that file started before, and keeps \code{rkl = true} and
\code{pcg\_tol} $= 10^{-11}$, so resuming, refining or filming it never switches the integrator in the middle of a ladder. A key
given in \code{[scheme]} always wins.

\subsection{The rotating frame}

In the gauge $\Phi = 0$ a stationary state with charge imbalance is not a fixed point of an explicit step: where its
phase turns at rate $\omega$, the step of that rotation, $|1 + i\omega h| > 1$, leaves a radial error $\propto
\omega^2 dt$, and KWT, which slows the radial mode and leaves the rotation alone, multiplies it by
$\sqrt{1+\gt^2|\Psi|^2}$. The amplitude in the conversion layers then depends on $dt$ at first order.

With \code{rotframe} each band is stepped in the gauge that turns with it: $\Psi = e^{i\hat\omega\tau}\tilde\Psi$, with
$\tau$ the time from the start of the step and $\hat\omega$ the tangential rate of the first stage at each node. The
band equation gains $-i\hat\omega\tilde\Psi$ (tangential, so the KWT factor leaves it alone), the links gain
$e^{i(\hat\omega_{i+1}-\hat\omega_i)\tau}$, and $\Psi^{n+1} = e^{i\hat\omega\,dt}\tilde\Psi$ turns the state back
exactly. A stationary state is then a fixed point of every stage for any $dt$. Any bounded $\hat\omega$ is an exact change
of gauge, so the two thresholds only decide where the frame helps: $\hat\omega = 0$ where $|\Psi|^2 < 10^{-3}$ (vortex
cores, the normal metal), and $|\hat\omega|\,dt \le 0.5$. Measured: with the frame the KWT state is the same to six
digits at $dt$, $dt/2$, $dt/4$ and with explicit substeps, for $\gt = 0$, 10 and 30; $\gt = 30$ diverged without it.

\textbf{Defaults.} \code{rotframe} and \code{rkl\_margin} = 1.1 are on when any band has KWT and off otherwise, so every
run without KWT keeps its bits.

\subsection{Cost}

Measured on one core, the two-band notched box ($40\times24\times6$ cells, $dx = 0.5$, $\kappa = 2$): without KWT
0.078~s per $\tGL$; with $\gt = 10$ and the defaults $\times2.2$. The two rotations of a link at each stage (its own and
the frame's) are linear in the fraction of the step, so their sum is computed once per band and every stage is one pass
with one $e^{ia}$; the rotating frame then costs nothing over the stages. The explicit step at the KWT bound would be ten
times shorter.
\code{dt\_scale = 1} halves the cost. On that box the delay of the first slip changed by $0.2\,\%$ without KWT (583
against 584 $\tGL$), and with KWT ($\gt = 10$) the silent state agreed to six digits and the delays at 0.74 and 0.75 were
the same: with the rotating frame an onset does not feel it. Flowing voltages were not checked; quote them at the default
step.

% ======================================================================
\section{Boundaries of the self-field}
\label{sec:bc}

The transport current enters as normal current through the two contact faces and its own Amp\`ere field is imposed on
the outer faces. Three choices (\code{selffield\_bc}):
\begin{description}
\item[\code{:slab}] the boundary of release 1.x: $h_y = \mp H_J$ on the whole $z$ faces. Amp\`ere then encloses the
  current of the sample and of the buffer, $J\,d\,(W + 2\,b\,dx - dx)$ for a buffer of $b$ cells: 4\,\% more than the
  nominal current on the $40\times24\times6$ box with one cell. It is the choice to reproduce release 1.x.
\item[\code{:bar}] (default) the field of the nominal current spread over the section of the sample, on every outer
  face, so the sample carries $J\,W\,d$.
\item[\code{:open}] the bar field plus the free-space field of every other current of the box, on all six faces. With
  $\delta\mathbf J = \kappa^2\nabla_h\times\mathbf h - \mathbf J_{\mathrm{bar}}$ (a closed current inside the box),
  $\mathbf A_\delta = G * \delta\mathbf J/\kappa^2$ and $\mathbf h_\delta = \nabla_h\times\mathbf A_\delta$ on the face
  plaquettes, where $G$ is the Green function of the seven-point lattice Laplacian ($-\Delta_h G = \delta$, $G(0) =
  0.252731\ldots$), so that the face sees the free continuation of the same discrete scheme. The convolution is
  Hockney's (zero padding to twice the box, FFT). $\mathbf h_\delta$ is recomputed every \code{open\_every} $\tGL$ (1 by
  default) and ramped in between; the iteration contracts (spectral radius at most 0.67, measured in 2D at $\kappa =
  2$). It is the field of the screening currents that a mirror face suppresses at small $\kappa$. It costs 15--20\,\% more
  per $\tGL$ than \code{:bar} on the box above ($3\,\%$ with KWT), keeps the $J\to-J$ symmetry at $H = 0$ and resumes
  exactly from a checkpoint.
\end{description}

% ======================================================================
\section{Measuring}

\subsection{The voltage and the verdict}

As in release 1.2: the voltage is the time derivative of the gauge-invariant phase difference
$\arg(\Psi_i^*U\Psi_j)$ summed along each measuring line between two planes \code{delta\_v} cells inside the contacts
(8 by default), unwrapped line by line, so that each $2\pi$ slip of a line is a vortex crossing. A rung integrated for
$\Delta t$ at voltage $V$ carries $|V|\Delta t/2\pi$ crossings per line: it is silent below one, dissipating at ten or
more, and undecided in between. The verdict reads band 1, so band 1 must be a band with a condensate of its own, not a
passive one. With noise the state never reaches a fixed point, so a rung closes silent when it held \code{t\_min\_GL\_res}
without a phase slip on any band; set that time to the waiting time the question needs (an escape time is a distribution, not a
number: for it, many seeds).

\subsection{Vortices}

\code{count\_vortices(G, Vx, Vy, hz, f, grid, k)} counts the vortices of one band on the $x$--$y$ plaquettes of plane $k$: the
gauge-invariant phase differences around a plaquette plus the flux through it are $2\pi$ times an integer, exactly on the
lattice, so the count is an integer and not a rounding (a vortex $+1$, an antivortex $-1$, along $+z$). Their sum over the plane
equals the winding of the phase around the edge of the sample (the discrete Stokes theorem; checked on a field-cooled lattice of
35 vortices, every plane and band), which is the fluxoid and not the flux$/2\pi$: the screening currents of the edges carry
field without vortices. Two limits: a plaquette counts only when its four corners conduct, and a band with almost no condensate
(a passive band above its own critical field) still gets the windings of the phase the Josephson coupling imposes on it, so a
count is read with $|\Psi_\mu|^2$ beside it. Lines along $x$ or $y$ are not in these maps. Count on the interior planes
($z_1 < k \le z_2$): the two face planes are copies of their neighbours, and the mask $f$ of a notch is written on the bottom one
but not on the top one (the masks are those of release 1.2, to the bit), so a count there sees the windings of the weak condensate
inside the notch. The dynamics never reads the face planes, and the tools of Sec.~\ref{sec:running} count and cut on interior
planes only.

\subsection{A current ladder}

Each current starts from the state the previous one left (the order of the ladder is part of the measurement: it gives
the hysteresis). A point is integrated in cycles of \code{lMax} steps until it closes:
\begin{description}
\item[\code{'F'}] a fixed point: $|\Psi_\mu|$ and $h$ change by less than the tolerance per cycle after the floor
  \code{t\_min\_GL} (400~$\tGL$ by default), and the second half carries fewer than \code{f\_margin} = 0.5 crossings
  per line (a slow flow near an onset can look fixed for a short cycle).
\item[\code{'V'}] dissipating: after the floor \code{t\_min\_GL\_res} (1584~$\tGL$), ten or more crossings per line, closed
  at \code{v\_margin} = 1.1 times that by the end-point estimate so that the slope the verdict uses lands above ten.
\item[\code{'N'}] a sample gone normal: a point that would close \code{'F'} with $\langle|\Psi_1|^2\rangle$ below \code{n\_frac} =
  0.02 of the largest of the \code{'F'} points before it. The normal state is a fixed point with no phase to cross; with large
  $u$ or $\gt$ an onset can be a collapse. \code{'N'} counts as dissipating.
\item[\code{'T'}] the heat model went above its cut.
\item[\code{'D'}] a point that closed \code{'F'} and crossed when continued (below).
\end{description}
\code{bracket\_Ic} gives the onset bracket of each polarity from the verdicts, and \code{run\_ladders} the pair and
$\eta$ with its range.

\subsection{The onset by the delay}\label{sec:delay}

Near a saddle-node onset the delay before the first phase slip, after a jump from a silent state, diverges as
\begin{equation}
  D(J) = \frac{a}{\sqrt{J - J_c}} + b .
\end{equation}
A window of fixed length therefore biases the onset up: the rung just above $J_c$ closes \code{'F'} on any finite window (in
the two-band diode at $\kappa = 2$, delays of 1600--2600~$\tGL$ against windows of 400). With
\code{Protocol(t\_confirm\_GL = ...)}, at the first \code{'V'} after an \code{'F'} the \code{'F'} rungs below it are continued at their
current, from copies of their closed states, up to that time in all (at most \code{confirm\_max} = 4): one that crosses
becomes \code{'D'} and the next one down is continued, until one stays silent. \code{delay\_bracket} gives the onset that way;
it is the one to quote. \code{crossing\_times(st, s, J; H, T)} measures delays directly from one silent state, and
\code{tools/delay\_fit.py} fits the law to two or more of them (with two, $b = 0$; it warns when the fit sits at the edge
of its search range).

% ======================================================================
\section{Optional terms}

Off by default and meant one at a time (the first three as in release 1.2):
\begin{description}
\item[Heat] a temperature field above the bath, $\varepsilon_\mu = a_\mu - \theta$ on every node,
  $\partial_t\theta = D_T\nabla^2\theta + \eta\,[\sum_\mu\Gamma_\mu|\partial_t\Psi_\mu|^2 + \sigma_n|\mathbf E|^2]
  - \gamma(\theta - T)$; the source is the dissipation of the functional, so a Meissner state does not heat. The Joule
  term is optional (\code{joule}), and \code{margin} cells next to the injection faces are left out of the source
  (contact heat, which leads would carry away). A point above the cut closes \code{'T'}.
\item[Drag] the mixed-gradient term $\nu_{\mu\nu}\mathrm{Re}[(\mathbf D\Psi_\mu)^*\cdot\mathbf D\Psi_\nu]$ of the
  functional: band $\mu$ feels $\nu_{\mu\nu}\mathbf D^2\Psi_\nu$ and the supercurrent carries the cross terms.
\item[Bath temperature] a uniform $T$ in $\varepsilon_\mu = a_\mu - T$ for every band.
\item[Two normal fluids] (\code{qp\_exchange} $=\nu$, two bands). With one normal current $\sigma_n\mathbf E$ the
  quasiparticles of every band share one potential: their interband exchange is infinitely fast. With a finite rate $\nu$
  each band has its own normal current $s_\mu\sigma_n(-\nabla\mu_\mu - \partial_t\mathbf A)$ ($s_\mu$ its share of
  $\sigma_n$, \code{qp\_split}, by default the dirty one $\propto g_\mu\Gamma_\mu/\beta_\mu$), and the bands exchange with
  a conductance $G = \chi_\sigma\nu$. In the gauge $\sum_\mu s_\mu\mu_\mu = 0$ Amp\`ere's law is unchanged and one field is
  new, $\delta = \mu_\pi - \mu_\sigma$:
  \begin{equation}
    \sigma_n\nabla^2\delta - G\Big(\frac{1}{s_\sigma} + \frac{1}{s_\pi}\Big)\delta = \frac{C_\pi}{s_\pi} - \frac{C_\sigma}{s_\sigma},
    \qquad C_\mu = \mathrm{Im}(\Psi_\mu^* R_\mu),
  \end{equation}
  $C_\mu$ the conversion of band $\mu$ (the divergence of its supercurrent plus the Josephson exchange), with $\delta = 0$ at
  the contacts and Neumann elsewhere, solved exactly in the eigenbases of the three one-dimensional operators. Each band
  then turns by $e^{-i\mu_\mu dt}$, tangential, so KWT does not act on it. It is the slaved branch imbalance with a finite
  exchange (its memory does not move a saddle-node onset); $\nu\to\infty$ is the step of one fluid to $O(1/\nu)$, and in a
  strip the $\sigma$ share of the converted current follows the one-dimensional two-fluid law. With $N$ bands
  \code{qp\_exchange} is the matrix of rates $\nu_{\mu\nu}$ (detailed balance $\chi_\mu\nu_{\mu\nu} = \chi_\nu\nu_{\nu\mu}$): with
  $L$ the Laplacian of the exchange graph and $S = \mathrm{diag}(s)$, $\sigma_n\nabla^2\mu - S^{-1}L\mu = S^{-1}C - \mathbb 1\sum C$,
  which the eigenvectors of $S^{-1/2}LS^{-1/2}$ split into $N$ Helmholtz problems of the same kind (with two bands, the $\delta$
  equation above). With the insulating buffer the quasiparticles live in the conductor: $\delta$ is solved on the bar of the
  sample and its contacts, with zero flux to the insulator (separable, exact; the strip then follows the one-dimensional law with
  the code's contact to 0.004, where the metal buffer, which conducts and exchanges around it, is 0.025 off). In a sculpted sample
  the links to the vacuum are cut and the equation is no longer separable: conjugate gradients in the inner product with weight
  $1/2$ on each Neumann face (where the mirror-ghost operator is symmetric), preconditioned by the separable solve of the bar and
  projected on the conductor, started cold every step (20--60 iterations). With the Joule heat the fluids add
  $\sigma_n\sum_\mu s_\mu|\nabla\mu_\mu|^2 + \sum_{\mu<\nu}G_{\mu\nu}(\mu_\mu-\mu_\nu)^2$ to $\sigma_n|\mathbf E|^2$ (the cross
  term vanishes in the gauge $\sum_\mu s_\mu\mu_\mu = 0$).
\item[Anisotropic bands] (\code{m\_axes}, or \code{m\_x}, \code{m\_y}, \code{m\_z} in a material file): a mass per
  axis in the band update, the supercurrent and the step bound ($4\sum_i g_i/dx_i^2$). An isotropic band keeps its bits;
  a band with $m_x = 4$ at $dx = 0.5$ is the isotropic one at $dx = 1$ to $10^{-9}$.
\item[Thermal noise] (\code{noise} $=\tilde T$, \code{noise\_seed}): a Langevin term
  $\langle\eta_\mu\eta_\mu^*\rangle = 2\Gamma_\mu\tilde T\delta(\mathbf r-\mathbf r')\delta(t-t')$ on every band, so that the
  discrete system is distributed as $e^{-F/\tilde T}$ (checked: mean, variance and correlation of decoupled nodes, and the
  spatial covariance against the linear theory of the lattice to 0.1--1\,\%). Seeded, saved in the checkpoints (resuming
  keeps the bits). With KWT the mobility of a node depends on $\Psi$, $M = (s/\Gamma)I - (\gt^2/\Gamma s)\psi\psi^{\mathsf T}$ with
  $s = \sqrt{1+\gt^2|\Psi|^2}$: the kick is scaled apart along $\Psi$ ($1/\Gamma s$) and across it ($s/\Gamma$), and the It\^o
  drift $\tilde T\nabla\cdot M/(2dV) = -\tilde T\gt^2(1+s^2)\psi/(2\Gamma s^3 dV)$ keeps $e^{-F/\tilde T}$ (checked: the same mean
  and variance as without KWT, the amplitude $s$ times slower and the phase $s$ times faster; without the drift the mean of
  $|\Psi|^2$ comes out 1.39 instead of 1). On a lattice the KWT phase modes are $s$ times faster and the noise excites them at
  the scale of the mesh: kicked after the stages they diverged (or, when finite, gave $\langle|\Psi|^2\rangle$ off by up to $-67\,\%$).
  So with noise a KWT band steps in log-polar form (its increments, the frame term and the combination of its stages act on the
  phase as rotations, and the frame rate is smoothed, an exact change of gauge), the phase kick enters the first RKC stage
  (SK-ROCK: Abdulle, Almuslimani and Vilmart, SIAM/ASA J.\ Uncertain.\ Quantif.\ 6, 937 (2018)), and the amplitude kick is
  polar, with the drift $\tilde T/(2\Gamma dV\rho s^3)$. With \code{noise\_field} the same bath acts on the field: Amp\`ere is
  $\sigma_l\partial_t\mathbf A = -\tfrac12\delta F/\delta\mathbf A$ per unit volume, so a link takes
  $\Delta A_l = \sqrt{\tilde T dt/(\sigma_l dV)}\,n$, added to the explicit part of the field step before the implicit solve. In a
  normal box the field energy of the free (transverse) modes is $\sum_m(\tilde T/2)/(1+\lambda_m dt/2)$ to 0.03\,\%, with
  $\sigma_n$ and with the insulating buffer: the implicit Euler of the field leaves the stiff mesh modes colder and the slow ones
  right to $O(\lambda dt)$.
\end{description}
None of them sets the step in practice (the heat diffusion bound, $0.15\,dx^2/D_T$, is 0.15 at $dx = 0.5$ and $D_T =
0.25$, against $dt \approx 0.02$), so an implicit treatment would gain nothing.

% ======================================================================
\section{Running}\label{sec:running}

\subsection{From a file}

A run is one TOML file:
\begin{verbatim}
[material]   file = "mgb2_model_II.toml"      (or kappa, sigma_n, T, gamma12 | gamma,
                                                and [[material.band]] tables inline)
[materials.NAME]  the same keys, for a shape region with material = "NAME"
             (+ transparency = { other = t })
[mesh]       nx, ny, nz, dx, buf, ghost_links
[geometry]   GEOM_MODE, geom_file, ...         (the shape keys of release 1.2)
[scheme]     the Stepper keywords: dt_scale, rkl (true | false | "rkc"), rotframe,
             rkl_margin, selffield_bc ("bar" | "open" | "slab"), open_every,
             sigma_vac, pcg_tol, ...
[protocol]   t_min_GL, t_min_GL_res, t_confirm_GL, ...
[heat]       D_T, gamma_bath, eta, joule, margin       (optional; drag goes in [scheme])
[run]        J = [...], H, out_dir, start
\end{verbatim}
and one command:
\begin{verbatim}
OPENBLAS_NUM_THREADS=1 julia --project=tdgl2 \
    tdgl2/scripts/run.jl RUN.toml [--dir OUT]
\end{verbatim}
Each polarity starts from the Meissner seed, or from what \code{start} names: \code{"field\_cooled"} (the sample normal at the
bath temperature \code{start\_T\_hot} = 1.5 for \code{start\_t\_hot} = 150~$\tGL$, so that the field fills it, then cooled from
small noise: the vortex state of a smooth sample at a field below its entry field, which the Meissner seed keeps out) or a state
file (\code{save\_state}; a path inside the folder of the run file is copied into the run folder, as the shape). It writes into the run folder the file it was given, \code{output/rungs\_J+.csv} and \code{rungs\_J-.csv} (one line per
rung: verdict, voltage, crossings, band currents and amplitudes, delays), \code{restart/} (JLD2 checkpoints, from which
a ladder resumes), and \code{output/summary.json}: the units, the code version and commit, the scheme as resolved (step,
the bound that set it, stages, \code{dt\_scale}, \code{rotframe}, \code{rkl\_margin}, the boundary), the protocol, the
regime numbers, and the brackets of the window and of the delay with $\eta$. \textbf{Stopping and resuming.} The checkpoint of every closed point holds, in one file written at once, the state, the
open boundary and the temperature field, the solution of the insulating solve, and all the ladder carries from one point
to the next (the convergence references, the reference of the 'N' verdict, the closed rungs waiting for the delay
confirmation). \code{STOP.txt} in the run folder stops a run after the current point; \code{--resume} continues it, and
the result is the one of the run without a stop, to the bit (the suite checks it, with the delay confirmation half done
and with the open boundary). A killed run loses only the point it was on. \code{--polarity +} and \code{--polarity -} run
the two polarities as two processes into the same folder; the second to end writes the summary, the same as one process
running both. A folder that already holds a ladder is not run over without \code{--resume}.

The examples are in \code{examples/}:
\code{wedge\_k25.toml} (minutes) and \code{kwt\_k2.toml} (about an hour; the onset with KWT, $\gt = 10$, rises to
$(0.73, 0.74]$ and $(0.75, 0.76]$ with $\eta = -1.34\,\%$); and one per option: \code{open\_k2},
\code{insulating\_k2}, \code{heat\_k25} (heat with Joule, a contact margin and drag), \code{three\_bands} and
\code{cable\_k2} (a sculpted nanowire in vacuum).

\subsection{From Julia}

\begin{verbatim}
using TDGL
md = two_band(xi_ratio=4.0, κ=2.0, gamma_kwt=10.0)
gr = Grid(40, 24, 6; dx=0.5)
mk = build_masks(GeomSpec(GEOM_MODE="WEDGE"), gr)
st = Stepper(md, gr, mk)                 # production scheme; rotframe on (KWT)
L  = ladder(st, initial_state(md, gr, mk, 0.03), [0.0, 0.5, 0.7, 0.74]; H=0.03,
            pr=Protocol(t_confirm_GL=3000.0, stop_after_V=1))
delay_bracket(L)
\end{verbatim}

\subsection{The terminal menu, campaigns and tools}

The terminal menu of release 1.2 (\code{julia/scripts/tdgl-tui.sh}) lists the runs of 2.0 that have an
\code{out\_dir}, with their points, and launches one from a run file (key \code{n}; with two CPUs, one process per
polarity), stops it after the current point (\code{s}) and resumes it (\code{u}). Its launch wizard (\code{l}) ends with
the choice of release: \code{2} writes the run file of release 2 (the two-band model of 1.2, KWT off, with the \code{:bar}
boundary of 2.0; \code{selffield\_bc = "slab"} is the one of 1.2) and launches it. In the details of a run, \code{f} draws its
figures, \code{m} films its vortex events and \code{w} opens the film.

\textbf{Figures and films} (Python 3 with NumPy and matplotlib; optional, the runs do not need it):
\begin{itemize}
\item \code{scripts/figures.jl RUNDIR}: the standard set into \code{RUNDIR/figures/}, with \code{README.txt} saying what each one
  is: the sheet of the ladders (\code{tools/ladder\_sheet.py}), $|\Psi_\mu|^2$ of every band on the three mid cuts with its
  vortices and the supercurrent lines, the share of the current of each band along the sample (\code{tools/state\_figures.py}),
  and \code{REPORT.md}, one page with what was run, every rung, the brackets and $\eta$ (\code{tools/run\_report.py}).
\item \code{scripts/film.jl RUNDIR [+|-] [--J J] [--events N]}: from the kept state of a rung (\code{keep\_rung\_states = true}), the
  frames around each vortex event only (the entry, the crossing and the exit; the waiting is skipped). It writes
  \code{film.html} (3D in a browser, turned with the mouse, with a slider; \code{tools/film\_html.py}), \code{film.gif}
  (\code{tools/film\_gif.py}) and the data as VTK files (\code{film.pvd}) for any VTK reader. A frame costs a few ms against a
  step of the solver, so the film does not slow the run. It steps the same equations from the kept state, so it shows the
  dynamics of that rung; with the insulating padding the solve starts cold and agrees with the ladder to its tolerance, not to
  the bit. With \code{--state FILE} it films any state file of the run's mesh instead (at \code{--J}, 0 by default: the cooling
  of a sample seeded by \code{start\_state}, a lattice forming). The face planes $y$ and $z$ are ghosts and are drawn as what they copy: in a box, their neighbour (a notch reads the
  same on both faces); in a sculpted sample, outside, because a link to a ghost carries no current as a link to vacuum does (a
  cylinder of a few cells draws round, not with a flat bulge where it meets the box). In a continuous flow every sample is an
  event, so \code{--every} must be well below the period of the flow, $2\pi/V$: the $\nu=1$ box of the diode at $J=0.93$, above its onset,
  slips every $1.05\,t_{\rm GL}$, and sampled every $2\,t_{\rm GL}$ its frames alias like a stroboscope (the film says so when the
  vortex map changes at most samples).
\item \code{scripts/refine.jl RUNDIR + | - [--step S]}: the bracket of a run refined from the state of its last silent rung (or a
  ladder that never flowed carried further up), into a subfolder that is a run folder itself; it reads the run and every
  refinement already made and prints the bracket and $\eta$. With \code{--confirm T} it confirms the bracket by the delay
  (Sec.~\ref{sec:delay}: the silent rungs below the flowing end followed $T$~$\tGL$; one that crosses becomes \code{D}).
\item \code{tools/collect\_runs.py DIR...}: one CSV with a row per run under some folders (options, brackets, $\eta$, version):
  a campaign as one table.
\end{itemize}

For campaigns, one single-threaded process per physical core (\code{OPENBLAS\_NUM\_THREADS=1}); on a laptop with hybrid
cores the total gain saturates near $2.5\times$ one process (memory bandwidth). The Slurm scripts of release 1.2
(\code{examples/}) serve with their command replaced by \code{tdgl2/scripts/run.jl RUN.toml --dir OUT}.

% ======================================================================
\section{Checks}

The suite (\code{tdgl2/test/run.jl}, a few minutes) checks:
\begin{itemize}
\item \textbf{Release 1.2.} The step, the masks of every catalogue shape, the checkpoints and a two-polarity ladder point
  by point, with the new terms off and \code{selffield\_bc = :slab}, against states written by release 1.2. To the bit
  with OpenBLAS's Haswell kernels and Julia code for that CPU (\code{TDGL\_TEST\_BITS=1 OPENBLAS\_CORETYPE=Haswell julia
  -C haswell}, any x86 machine with AVX2); to $10^{-10}$ elsewhere, where the dense products of the separable solve and
  the SIMD sums round differently (measured: $10^{-15}$ on the states, $10^{-12}$ on ladder voltages, which are time
  averages). Verdicts and counts are exact either way.
\item \textbf{Energy.} The discrete balance $dF/dt = -D + W$ on a two-band box relaxing from a random state, where $W$
  is the work of the supercurrent on the links between two face nodes (about 0.5\,\% of $D$ there; with it the residual
  is $O(dt)$). With those links off, $W = 0$ and $F$ decreases at every step: an exact gradient flow.
\item \textbf{The rotating frame.} The KWT state at $dt$ and $dt/2$ and with substeps, to $10^{-5}$.
\item \textbf{Three bands.} Equal repulsive couplings frustrate the phases: the ground state breaks time reversal with
  phases $0$, $\pm2\pi/3$ and $|\Psi|^2 = 1 + |\gamma|$, the analytic values, to $10^{-6}$.
\item \textbf{Symmetry.} A sample with the mirror $y\to-y$ at $H\neq0$: the state at $+J$ is the mirror of the state at
  $-J$ to rounding ($10^{-15}$), with each boundary, so $\eta = 0$.
\item \textbf{The open boundary.} The lattice Green function, $\mathbf h_\delta = 0$ without currents, a closed source,
  the nominal current through the sample to $2\times10^{-3}$, exact resumption, the $J\to-J$ symmetry at $H = 0$.
\item \textbf{The protocol.} The verdicts \code{'N'} and \code{'D'}, \code{crossing\_times}, \code{delay\_bracket}, the regime numbers
  (one band against the conversion law).
\end{itemize}

\textbf{Against pyTDGL} (version 0.9; one band in 2D, the one-band equation of this code with $\Gamma = u$, $\sigma_n =
1$, $\kappa$ large; pyTDGL's current unit is four times the GL unit at its interface, and it solves in the GL unit inside):
the depairing current of a narrow strip, $0.385$; the conversion length with KWT for $\gt = 0$--30, $0.2\,\%$ between the
codes and $0.1\,\%$ from $\lambda_c^2 = \sqrt{1+\gt^2|\Psi|^2}/(u|\Psi|^2)$; the onset of a notched strip, independent of
$\gt$ in both codes, which with the mesh approaches pyTDGL's value ($8\,\%$ below it at $dx = 0.5\,\xi$, $4\,\%$ at
$0.25\,\xi$: the $O(dx)$ error at the notch tip). The flowing voltage has not been compared.

% ======================================================================
\section{Limitations}

\begin{itemize}
\item \textbf{Absolute onsets carry an $O(dx)$ error} of 5--12\,\% at $dx = 0.5\,\xi$, at any $\kappa$: the tip of a
  notch is resolved at the scale of a cell, and at small $\kappa$ the face rows of the sample, Neumann copies without an
  equation of their own, carry a large share of the screening supercurrent (51\,\% at $\kappa = 2$, $dx = 0.5$; 30\,\%
  at 0.25). Ratios between polarities ($\eta$) and trends do not feel it; quote absolute onsets with their mesh. The links
  that join two ghost nodes ($\tau = 1$, as in 1.x) are not variational: \code{ghost\_links = false} in \code{[mesh]} turns
  them off (the exact gradient flow, the surface half a cell in). With materials it is needed where a current crosses an
  interface that reaches the faces: the face links copy each side's $|\Psi|$ and their current jumps there, a spurious source
  of normal current (2.6\,\% of the supercurrent in a thin S/S$'$/S strip).
\item \textbf{With KWT the conversion at the contacts is long}: keep $L/\lambda_c \gtrsim 10$.
\item \textbf{The silent state is driven.} The conversion layers carry a steady field even when the order parameter is
  static, so an onset depends on the relaxation rates and on the contact, not only on the free energy. Quote onsets for
  the rates, the KWT parameters and the boundary used.
\item \textbf{The charge imbalance is slaved}: KWT per band is the branch (Schmid--Sch\"on) imbalance in its slaved form,
  and \code{qp\_exchange} adds a finite interband exchange; its memory and the scalar potential of 1.2 are not in. The normal
  fluids depend on the share of $\sigma_n$ between the bands, a model choice.
\item \textbf{The insulating buffer cost 8--10 times per $\tGL$} what a metal buffer costs at the tolerance $10^{-11}$ (since 26
  October 2026, with $10^{-8}$, the extrapolated start and the single-precision preconditioner, about three times: 4.5 ms per step on
  the nanowire box against 1.5 ms without vacuum) (two cells of padding around the
  notched box: $\times8$; the nanowire box, mostly vacuum: $\times10$, 23 iterations of the solve per step). The iterations are
  set by the smooth solenoidal fields of the vacuum (the middle of the preconditioned spectrum, $\kappa_{\mathrm{ef}}\approx4.5$,
  nearly independent of $\sigma_{\mathrm{vac}}$), which no cheap separable preconditioner removes, so a thicker padding costs more.
  The lever was the tolerance: \code{pcg\_tol} $= 10^{-8}$, the default, is $-41\,\%$ per step on the nanowire against $10^{-11}$
  and $-45\,\%$ on a ladder of the notched nanowire, with the same bracket and the flowing voltages to 9 digits. The padding
  thickness: one cell against two gives the same brackets there and flowing voltages 2--7\,\% higher. A sculpted sample needs it. The onset does not depend on
  $\sigma_{\mathrm{vac}}$ between 0.003 and 0.1 (the notched nanowire: the same brackets), the resistive voltage does ($-16\,\%$
  with 0.1): $\sigma_{\mathrm{vac}} = 0.1$ is $1.5$ times faster for onsets and $\eta$, 0.01 is the value for voltages.
\item \textbf{Anisotropy} is a mass per axis and band; $\sigma_n$ and drag stay isotropic.
\item \textbf{Noise and the heat model} are not combined: the heat source is the dissipation of the step, which with noise
  counts the Langevin kicks ($\propto\tilde T/dt$), so a sample in equilibrium would heat by itself.
\item \textbf{Noise with KWT} has an integrator error on the stiffest phase modes: $\langle|\Psi|^2\rangle$ within 1.2\,\% of the
  value without KWT at $\tilde T = 10^{-3}$, but $+5$--$10\,\%$ for the stiff $\pi$ band at $\tilde T = 10^{-2}$ at the production
  step ($-7\,\%$ with Euler substeps), which goes away only when the step resolves those modes. Check it with \code{dt\_scale}.
  On RKL stages it is refused.
\item \textbf{The implicit solve needs $dx = dy = dz$.}
\item \textbf{No GPU}: the loops are per node and per link with no reductions, so a port is mechanical; it has not been done.
\item \textbf{One single-threaded process per physical core} for campaigns. Within one run, threads give the same bits
  but no speed on a laptop (measured: $\times0.8$--$1.06$ with 2--8 threads).
\end{itemize}

\end{document}
